\documentclass[zihao=-4, a4paper]{ctexart}
\usepackage[margin=2.5cm]{geometry}
\usepackage{amsmath}
\usepackage{booktabs}
\usepackage{graphicx}
\usepackage{caption}
\usepackage{setspace}

\onehalfspacing

% ===== 课程信息（按需修改）=====
\newcommand{\coursename}{课程名称：数据结构}
\newcommand{\labtitle}{实验一：线性表的应用}
\newcommand{\studentname}{姓名：张三 \quad 学号：2022010101}
\newcommand{\labdate}{日期：\today}

\begin{document}

\begin{center}
  {\zihao{3}\heiti \labtitle}\\[0.6em]
  {\zihao{-4} \coursename}\\[0.3em]
  {\zihao{-4} \studentname}\\[0.3em]
  {\zihao{-4} \labdate}
\end{center}

\section{实验目的}
简述本次实验要掌握的知识点与能力，例如：掌握线性表的顺序存储与链式存储
实现，理解两种结构在插入删除操作上的复杂度差异。

\section{实验原理}
简述实验涉及的理论基础。例如顺序表按位置访问时间为 $O(1)$，
插入删除平均移动一半元素为 $O(n)$；单链表插入删除只需修改指针，
但查找为 $O(n)$。

\section{实验内容与步骤}
\begin{enumerate}
  \item 实现顺序表的基本操作（初始化、插入、删除、按位查找）；
  \item 实现单链表的基本操作；
  \item 分别在两种结构上执行 $10^4$ 次随机插入/删除，记录耗时；
  \item 对比分析实验数据。
\end{enumerate}

\section{实验结果}
\begin{table}[h]
  \centering
  \caption{两种线性表操作耗时对比（ms）。}
  \begin{tabular}{lcc}
    \toprule
    操作 & 顺序表 & 单链表 \\
    \midrule
    随机插入 & 12.4 & 3.1 \\
    随机删除 & 11.8 & 3.0 \\
    按位查找 & 0.1 & 8.9 \\
    \bottomrule
  \end{tabular}
\end{table}

\section{结果分析与结论}
结合数据说明结构选择与场景的关系，总结实验收获，也可以提出改进方向
（如双向链表、静态链表等）。

\end{document}
